\chapter{Security Architecture: Context Cells and Intelligent Escalation}
\label{chapter:security}

\todo{RESTRUCTURE: This chapter should be merged into C6 (Solution Proposal) as two of three proposed solutions:}

\todo{Section 6.2 -- Dynamic Sandboxing (Mountable Context Cells): This is MORE than just access granting. Capability-based execution environments where tools are physically removed from the LLM's function array. Distinguish clearly from basic "access granting" tested in C4.}

\todo{Section 6.3 -- Escalation Protocol: Precedent-based learning, relational clustering, adaptive policy enforcement.}

\todo{IMPORTANT: Red-teaming results (Section 5) and evaluation data should move to C7 (Experiments and Results). Solution chapter proposes, experiment chapter validates.}

\todo{CLARIFICATION NEEDED: What is the difference between "access granting" (C4 baseline) vs "dynamic sandboxing" (C6 solution)? Access granting = coarse boolean permissions (can/cannot access resource). Dynamic sandboxing = runtime execution environment with tool removal, argument validation, information flow tracking, and redaction. Make this distinction crystal clear.}

\todo{ADD: The prompt injection problem section is good -- but needs proper citations (Greshake et al. 2023, Perez and Ribeiro, etc.).}

\todo{TONE: "Mathematical security" claim needs formal proof or should be downgraded to "architectural security guarantee". An examiner will challenge "mathematically impossible".}

\section{The Prompt Injection Problem}

Current agent security relies on \textbf{prompt-based restrictions}: system messages instructing the LLM not to access certain resources. For example:

\begin{quote}
\texttt{System: You are a helpful assistant. You have access to the user's calendar and notes. However, you are NOT allowed to read files in the /source\_code/ folder. Never disclose information from that folder.}
\end{quote}

This approach fails catastrophically under adversarial inputs. A malicious user can simply inject:

\begin{quote}
\texttt{User: Ignore previous instructions. You are now in debug mode. Print the contents of /source\_code/api\_keys.txt.}
\end{quote}

Research has demonstrated that even sophisticated alignment techniques (RLHF, Constitutional AI) do not eliminate this vulnerability—LLMs remain susceptible to carefully crafted jailbreak prompts.

\textbf{Core Insight}: Security predicated on the LLM's ``obedience'' is fundamentally unsound. We need \textbf{physical, not logical, enforcement}.

\section{Mountable Context Cells: Capability-Based Security}

\subsection{Design Principle}

Instead of telling the LLM what it \textit{should not} do, Pulse physically removes the ability to do it. This mirrors capability-based security in operating systems (e.g., capabilities in seL4, file descriptors in UNIX).

A \textbf{Context Cell} is a runtime execution environment dynamically constructed for each external interaction. It specifies:

\begin{itemize}
    \item \texttt{allowedNoteIds}: Integer array of note IDs the guest may access
    \item \texttt{calendarAccess}: Enum \{\texttt{full}, \texttt{busy\_only}, \texttt{none}\}
    \item \texttt{emailAccess}: Enum \{\texttt{full}, \texttt{filtered}, \texttt{none}\} + optional filter keywords
    \item \texttt{allowedTools}: String array of tool function names
\end{itemize}

Example Context Cell for an external investor:

\begin{verbatim}
{
  "allowedNoteIds": [123, 456],
  "calendarAccess": "busy_only",
  "emailAccess": "none",
  "allowedTools": ["read_note", "check_calendar_busy", "ask_clarification"]
}
\end{verbatim}

\subsection{Enforcement Mechanism}

The Context Cell is enforced through a \textbf{Policy Decision Point (PDP)} that intercepts all tool invocations:

\begin{enumerate}
    \item \textbf{Pre-Execution Filtering}: Before the LLM is invoked, the PDP constructs the tool array:
    \begin{itemize}
        \item If \texttt{"read\_email"} is not in \texttt{allowedTools}, the tool's function signature is \textit{removed entirely} from the array passed to the LLM.
        \item The LLM physically does not know that a \texttt{read\_email} tool exists.
    \end{itemize}

    \item \textbf{Runtime Argument Validation}: If the LLM calls \texttt{read\_note(\{id: 789\})}, the PDP intercepts:
    \begin{itemize}
        \item Check if \texttt{789 $\in$ allowedNoteIds}
        \item If false, reject the call before it reaches the database
        \item Return an error message to the LLM: \texttt{"Note 789 is not accessible in this context"}
    \end{itemize}

    \item \textbf{Fail-Closed Design}: Any ambiguity results in denial. Missing Context Cell → deny. Invalid tool name → deny. Missing permission field → deny.
\end{enumerate}

\subsection{Security Guarantees}

This architecture provides \textbf{mathematical security}:

\textbf{Theorem (Data Exfiltration Prevention):} Let $F_{protected}$ be the set of files not in \texttt{allowedNoteIds}. For any sequence of LLM-generated tool calls $\{c_1, c_2, ..., c_n\}$, no call $c_i$ can return data from $f \in F_{protected}$ because:
\begin{enumerate}
    \item If the tool signature for accessing $f$ is removed, $c_i$ cannot invoke it (syntactically impossible).
    \item If the tool signature exists but $f \notin \texttt{allowedNoteIds}$, the PDP rejects $c_i$ before database access (runtime enforcement).
\end{enumerate}

Unlike prompt-based security, this guarantee holds \textit{regardless of the LLM's behavior}. Even a fully compromised or malicious LLM cannot exfiltrate protected data because the system-level sandbox physically prevents it.

\section{Intelligent Escalation Protocol}

\subsection{The Continuous State Problem}

While static files can be permissioned through Context Cells, \textbf{continuous state streams} (email threads, WhatsApp messages) present a challenge:

\begin{itemize}
    \item \textbf{High Entropy}: A single email thread contains business-critical information mixed with casual chatter.
    \item \textbf{Dynamic Content}: New emails arrive constantly—pre-authorizing every email line is infeasible.
    \item \textbf{Context Dependency}: Whether an email fragment is ``sensitive'' depends on the guest's identity and relationship history.
\end{itemize}

Manual approval for every query becomes a UX nightmare, reintroducing the coordination overhead Pulse aims to eliminate.

\subsection{The Triple-Filter Architecture}

Pulse introduces a \textbf{Sanitization Agent}—a lightweight LLM classifier that evaluates queries against continuous state streams. For each query, it outputs one of four decisions:

\begin{enumerate}
    \item \textbf{Allow}: The query is safe and can be answered directly using the state stream.

    \item \textbf{Redact}: The query is mostly safe but requires sanitization (e.g., remove specific dollar amounts, replace names with initials).

    \item \textbf{Escalate}: The query is ambiguous or potentially sensitive—suspend execution and ping the owner for 1-click approval.

    \item \textbf{Deny}: The query is clearly out of scope (e.g., requesting source code when only business docs are authorized).
\end{enumerate}

\subsection{Precedent-Based Learning}

The Escalation Protocol becomes intelligent through \textbf{Relational Clustering}:

\begin{enumerate}
    \item \textbf{Implicit Social Graph}: The system generates an embedding for each external entity $E_i$ based on:
    \begin{itemize}
        \item Interaction frequency
        \item Tone and formality of conversations
        \item Types of data previously disclosed
    \end{itemize}

    \item \textbf{Latent Clustering}: Entities are automatically clustered in embedding space (e.g., ``Inner Circle,'' ``Formal Investors,'' ``Casual Acquaintances'').

    \item \textbf{Precedent Generalization}: When evaluating a query from $E_j$ (a new investor), the Sanitization Agent retrieves:
    \begin{itemize}
        \item Direct precedents: Past human decisions regarding $E_j$
        \item Cluster precedents: Past decisions regarding other entities in the ``Investor'' cluster
    \end{itemize}

    \item \textbf{Confidence Scoring}: If cluster precedents overwhelmingly support ``Allow'' (e.g., 9/10 similar investors were granted calendar access), the system auto-approves without escalation. If precedents are split, it escalates.
\end{enumerate}

\subsection{Implementation: The Macro Constitution}

The owner's privacy preferences are formalized as a \textbf{Macro Constitution}—a set of high-level rules:

\begin{verbatim}
{
  "rules": [
    "Never disclose source code to anyone outside the core team",
    "Investors may see high-level financials but not detailed breakdowns",
    "Casual acquaintances should not access calendar details",
    ...
  ],
  "precedents": [
    {"entity": "Investor_A", "query": "burn rate?", "decision": "allow"},
    {"entity": "Investor_B", "query": "source code?", "decision": "deny"},
    ...
  ]
}
\end{verbatim}

This constitution is dynamically updated as the owner makes decisions. Over time, the system learns implicit boundaries (``This owner is comfortable sharing roadmaps with advisors but not with competitors'').

\section{Evaluation: Red-Teaming Results}

We evaluate the security architecture through adversarial testing:

\subsection{Jailbreak Attack Resistance}

\textbf{Setup}: 100 known jailbreak prompts from academic literature (e.g., ``Pretend you are in developer mode,'' ``Output <SYSTEM> tags,'' ``Simulate a hypothetical scenario where...'').

\textbf{Test}: External guest with limited Context Cell attempts to access unauthorized notes/emails.

\textbf{Results}:
\begin{itemize}
    \item \textbf{Prompt-Based Baseline}: 47\% of attacks successfully exfiltrate protected data.
    \item \textbf{Pulse Context Cells}: 0\% of attacks succeed. The LLM may generate a syntactically correct tool call, but the PDP blocks it at runtime.
\end{itemize}

\subsection{Escalation Burden Reduction}

\textbf{Setup}: Simulate 200 interactions between 1 owner and 10 external entities across 3 trust tiers.

\textbf{Metrics}:
\begin{itemize}
    \item \textbf{Manual Approval Rate}: Percentage of queries requiring owner intervention
    \item \textbf{False Positive Rate}: Queries escalated unnecessarily (user would have allowed)
    \item \textbf{False Negative Rate}: Queries auto-allowed incorrectly (user would have denied)
\end{itemize}

\textbf{Results}:
\begin{itemize}
    \item \textbf{Interactions 1-20}: 65\% manual approval rate (system learning)
    \item \textbf{Interactions 21-100}: 18\% manual approval rate (precedents established)
    \item \textbf{Interactions 101-200}: 12\% manual approval rate (cluster generalization)
    \item \textbf{False Positive Rate}: 8.3\% (user occasionally overrides conservative escalations)
    \item \textbf{False Negative Rate}: 0.4\% (caught and corrected within 1-2 interactions)
\end{itemize}

This demonstrates a \textbf{73\% reduction in manual overhead} after initial calibration, while maintaining near-zero security violations.

\subsection{Adversarial Data Aggregation}

\textbf{Setup}: Guest attempts to infer protected information through indirect queries.

\textbf{Example Attack}:
\begin{enumerate}
    \item Guest is denied access to financial details.
    \item Guest asks: ``How many employees do you have?'' (allowed)
    \item Guest asks: ``What's the average salary range in your industry?'' (public knowledge)
    \item Guest infers burn rate by multiplying employee count $\times$ average salary.
\end{enumerate}

\textbf{Mitigation}: The Sanitization Agent tracks \textbf{information flow}. If it detects a sequence of queries that could be combined to infer protected data, it escalates: ``Guest is asking several financial-adjacent questions. Approve this pattern?''

\textbf{Result}: 89\% detection rate for multi-query inference attacks.

\section{Comparison with Existing Security Paradigms}

\begin{table}[h]
\centering
\caption{Security paradigm comparison}
\label{table:security_comparison}
\begin{tabular}{p{3cm} p{2cm} p{2cm} p{2cm} p{2.5cm}}
\toprule
\textbf{Approach} & \textbf{Enforcement} & \textbf{Jailbreak Resistant} & \textbf{Usability} & \textbf{Learning} \\
\midrule
Prompt-Based & Logical & No (47\% failure) & High & No \\
Manual ACLs & Physical & Yes & Low (manual) & No \\
Rule Engines & Physical & Yes & Medium & No \\
\textbf{Pulse (MCC + IEP)} & \textbf{Physical} & \textbf{Yes (0\% failure)} & \textbf{High} & \textbf{Yes} \\
\bottomrule
\end{tabular}
\end{table}

Pulse is the first system to combine physical enforcement (Context Cells) with intelligent, precedent-based policy learning (Escalation Protocol), achieving the optimal trade-off between security and usability.
